%%%%%%%%%%%%%%%%%%%%%%%%%%%%%%%%%%%%%%%%%%%%%%%%%%%%%%%%%%%%%%%%%%%%%%%%
%%%%%%%%%%%%%%%%%%%%%% Simple LaTeX CV Template %%%%%%%%%%%%%%%%%%%%%%%%
%%%%%%%%%%%%%%%%%%%%%%%%%%%%%%%%%%%%%%%%%%%%%%%%%%%%%%%%%%%%%%%%%%%%%%%%

%%%%%%%%%%%%%%%%%%%%%%%%%%%%%%%%%%%%%%%%%%%%%%%%%%%%%%%%%%%%%%%%%%%%%%%%
%% NOTE: If you find that it says                                     %%
%%                                                                    %%
%%                           1 of ??                                  %%
%%                                                                    %%
%% at the bottom of your first page, this means that the AUX file     %%
%% was not available when you ran LaTeX on this source. Simply RERUN  %%
%% LaTeX to get the ``??'' replaced with the number of the last page  %%
%% of the document. The AUX file will be generated on the first run   %%
%% of LaTeX and used on the second run to fill in all of the          %%
%% references.                                                        %%
%%%%%%%%%%%%%%%%%%%%%%%%%%%%%%%%%%%%%%%%%%%%%%%%%%%%%%%%%%%%%%%%%%%%%%%%

%%%%%%%%%%%%%%%%%%%%%%%%%%%% Document Setup %%%%%%%%%%%%%%%%%%%%%%%%%%%%

% Don't like 10pt? Try 11pt or 12pt
\documentclass[10pt]{article}

% This is a helpful package that puts math inside length specifications
\usepackage{calc}
\usepackage{amsmath}
\usepackage{array}


% Simpler bibsection for CV sections
% (thanks to natbib for inspiration)
\makeatletter
\newlength{\bibhang}
\setlength{\bibhang}{1em}
\newlength{\bibsep}
 {\@listi \global\bibsep\itemsep \global\advance\bibsep by\parsep}
\newenvironment{bibsection}%
        {\vspace{-\baselineskip}\begin{list}{}{%
       \setlength{\leftmargin}{\bibhang}%
       \setlength{\itemindent}{-\leftmargin}%
       \setlength{\itemsep}{\bibsep}%
       \setlength{\parsep}{\z@}%
        \setlength{\partopsep}{0pt}%
        \setlength{\topsep}{0pt}}}
        {\end{list}\vspace{-.6\baselineskip}}
\makeatother

% Layout: Puts the section titles on left side of page
\reversemarginpar

%
%         PAPER SIZE, PAGE NUMBER, AND DOCUMENT LAYOUT NOTES:
%
% The next \usepackage line changes the layout for CV style section
% headings as marginal notes. It also sets up the paper size as either
% letter or A4. By default, letter was used. If A4 paper is desired,
% comment out the letterpaper lines and uncomment the a4paper lines.
%
% As you can see, the margin widths and section title widths can be
% easily adjusted.
%
% ALSO: Notice that the includefoot option can be commented OUT in order
% to put the PAGE NUMBER *IN* the bottom margin. This will make the
% effective text area larger.
%
% IF YOU WISH TO REMOVE THE ``of LASTPAGE'' next to each page number,
% see the note about the +LP and -LP lines below. Comment out the +LP
% and uncomment the -LP.
%
% IF YOU WISH TO REMOVE PAGE NUMBERS, be sure that the includefoot line
% is uncommented and ALSO uncomment the \pagestyle{empty} a few lines
% below.
%

%% Use these lines for letter-sized paper
\usepackage[paper=letterpaper,
            includefoot, % Uncomment to put page number above margin
            marginparwidth=1.2in,     % Length of section titles
            marginparsep=.05in,       % Space between titles and text
            margin=0.5in,               % 1 inch margins
            includemp]{geometry}

%% Use these lines for A4-sized paper
%\usepackage[paper=a4paper,
%            %includefoot, % Uncomment to put page number above margin
%            marginparwidth=30.5mm,    % Length of section titles
%            marginparsep=1.5mm,       % Space between titles and text
%            margin=25mm,              % 25mm margins
%            includemp]{geometry}

%% More layout: Get rid of indenting throughout entire document
\setlength{\parindent}{0in}

%% This gives us fun enumeration environments. compactitem will be nice.
\usepackage{paralist}

%% Reference the last page in the page number
%
% NOTE: comment the +LP line and uncomment the -LP line to have page
%       numbers without the ``of ##'' last page reference)
%
% NOTE: uncomment the \pagestyle{empty} line to get rid of all page
%       numbers (make sure includefoot is commented out above)
%
\usepackage{fancyhdr,lastpage}
\pagestyle{fancy}
\pagestyle{empty}      % Uncomment this to get rid of page numbers
\fancyhf{}\renewcommand{\headrulewidth}{0pt}
\fancyfootoffset{\marginparsep+\marginparwidth}
\newlength{\footpageshift}
\setlength{\footpageshift}
          {0.5\textwidth+0.5\marginparsep+0.5\marginparwidth-2in}
\lfoot{\hspace{\footpageshift}%
       \parbox{4in}{\, \hfill %
                    \arabic{page} of \protect\pageref*{LastPage} % +LP
%                    \arabic{page}                               % -LP
                    \hfill \,}}

% Finally, give us PDF bookmarks
\usepackage{color,hyperref}
\definecolor{darkblue}{rgb}{0.0,0.0,0.5}
\hypersetup{colorlinks,breaklinks,
            linkcolor=darkblue,urlcolor=darkblue,
            anchorcolor=darkblue,citecolor=darkblue}

%%%%%%%%%%%%%%%%%%%%%%%% End Document Setup %%%%%%%%%%%%%%%%%%%%%%%%%%%%


%%%%%%%%%%%%%%%%%%%%%%%%%%% Helper Commands %%%%%%%%%%%%%%%%%%%%%%%%%%%%

% The title (name) with a horizontal rule under it
%
% Usage: \makeheading{name}
%
% Place at top of document. It should be the first thing.
\newcommand{\makeheading}[1]%
        {\hspace*{-\marginparsep minus \marginparwidth}%
         \begin{minipage}[t]{\textwidth+\marginparwidth+\marginparsep}%
                {\large \bfseries #1}\\[-0.15\baselineskip]%
                 \rule{\columnwidth}{1pt}%
         \end{minipage}}

% The section headings
%
% Usage: \section{section name}
%
% Follow this section IMMEDIATELY with the first line of the section
% text. Do not put whitespace in between. That is, do this:
%
%       \section{My Information}
%       Here is my information.
%
% and NOT this:
%
%       \section{My Information}
%
%       Here is my information.
%
% Otherwise the top of the section header will not line up with the top
% of the section. Of course, using a single comment character (%) on
% empty lines allows for the function of the first example with the
% readability of the second example.
\renewcommand{\section}[2]%
        {\pagebreak[2]\vspace{1.3\baselineskip}%
         \phantomsection\addcontentsline{toc}{section}{#1}%
         \hspace{0in}%
         \marginpar{
         \raggedright \scshape #1}#2}

% An itemize-style list with lots of space between items
\newenvironment{outerlist}[1][\enskip\textbullet]%
        {\begin{itemize}[#1]}{\end{itemize}%
         \vspace{-.6\baselineskip}}

% An environment IDENTICAL to outerlist that has better pre-list spacing
% when used as the first thing in a \section
\newenvironment{lonelist}[1][\enskip\textbullet]%
        {\vspace{-\baselineskip}\begin{list}{#1}{%
        \setlength{\partopsep}{0pt}%
        \setlength{\topsep}{0pt}}}
        {\end{list}\vspace{-.6\baselineskip}}

% An itemize-style list with little space between items
\newenvironment{innerlist}[1][\enskip\textbullet]%
        {\begin{compactitem}[#1]}{\end{compactitem}}

% An environment IDENTICAL to innerlist that has better pre-list spacing
% when used as the first thing in a \section
\newenvironment{loneinnerlist}[1][\enskip\textbullet]%
        {\vspace{-\baselineskip}\begin{compactitem}[#1]}
        {\end{compactitem}\vspace{-.6\baselineskip}}

% To add some paragraph space between lines.
% This also tells LaTeX to preferably break a page on one of these gaps
% if there is a needed pagebreak nearby.
\newcommand{\blankline}{\quad\pagebreak[2]}

% Uses hyperref to link DOI
\newcommand\doilink[1]{\href{http://dx.doi.org/#1}{#1}}
\newcommand\doi[1]{doi:\doilink{#1}}


%%%%%%%%%%%%%%%%%%%%%%%% End Helper Commands %%%%%%%%%%%%%%%%%%%%%%%%%%%

%%%%%%%%%%%%%%%%%%%%%%%%% Begin CV Document %%%%%%%%%%%%%%%%%%%%%%%%%%%%

\begin{document}
\makeheading{ \hfill{Russell Cohen}}

\vspace{-5mm}

\section{Contact Information}
%
% NOTE: Mind where the & separators and \\ breaks are in the following
%       table.
%
% ALSO: \rcollength is the width of the right column of the table
%       (adjust it to your liking; default is 1.85in).
%
\newlength{\rcollength}\setlength{\rcollength}{1.85in}%
%
%\begin{tabular}[t]{@{}p{\textwidth/3}p{\textwidth/3}p{\textwidth/3}}
%\textit{email:} &\textit{US address:} & \textit{Permanent address:}  \\
 %\href{mailto:bjoveski@mit.edu}{bjoveski@mit.edu} & 28 The Fenway & ``Partizanski Odredi" \\
%\textit{cell phone:} & Boston, MA & br. $72^a/48$ \\
 %(617) 894-8789  & 02215        &  Skopje, Macedonia \\
%\end{tabular}
%\begin{center}
\begin{tabular}{>{\centering\arraybackslash}m{\textwidth}}
\\
\href{mailto:rcoh@rcoh.me}{rcoh@rcoh.me}
\end{tabular}
%\end{center}


%\section{Objective}
%
%Obtaining an internship position at Google's BOLD Practicum for the summer 2011

\section{Work Experience}
\href{https://github.com/awslabs/aws-sdk-rust}{\textbf{AWS}} \hfill{Senior Software Engineer, July 2020 | Present}
 \begin{outerlist}
    \item[] \begin{innerlist}[-]
        \item Lead the team supporting Rust adoption across Amazon. Find great ideas from teams
        around the company and distill them into core primitives that become de facto solutions.
        \item Created \href{https://github.com/awslabs/metrique}{metrique}, a low-overhead, wide-event
        metrics library. It replaced effectively all hand-written metrics libraries at Amazon.
        \item Created \href{https://github.com/dial9-rs/dial9}{dial9}, a flight recorder and profiler
        that records Tokio runtime events, application traces, and kernel events on a single timeline.
        Runs in production at Amazon, Cloudflare, OpenAI, and many others.
        \item Founded the \href{https://github.com/awslabs/aws-sdk-rust}{AWS SDK for Rust} team and
        led it to GA. At 20--30$\mu$s of overhead per request, it is among the fastest AWS SDKs.
        \item Built \href{https://github.com/awslabs/smithy-rs}{smithy-rs}, the code generator behind
        the Rust SDK, dozens of internal AWS services, and several external customers.
    \end{innerlist}
\end{outerlist}

\blankline

\href{http://rcoh.me}{\textbf{Self Employed}}, Redwood City, CA \hfill{Freelance Software Engineer, February 2016 | 2020}
 \begin{outerlist}
    \item[] \begin{innerlist}[-]
        \item CTO/senior engineer for hire: consulted on architecture and scalability, and joined
        startups for 1-3 months to finish high-value projects.
        \item Wrote \href{https://github.com/rcoh/angle-grinder}{angle-grinder}, a popular Rust CLI app for
        log analytics, and \href{https://crossword.tools}{Crossword.Tools}, a crossword construction tool
        with a Rust fill algorithm compiled to WASM.
    \end{innerlist}
\end{outerlist}

\blankline

\href{http://www.sumologic.com/}{\textbf{Sumo Logic}}, Redwood City, CA \hfill{Software Engineer Lead, February 2014  | February 2016}
 \begin{outerlist}
	\item[] \begin{innerlist}[-]
		\item Led the search performance team; developed the optimizer for the Sumo Logic query language.
		\item Authored \href{https://github.com/SumoLogic/elasticsearch-client}{open source ElasticSearch library}
	    \end{innerlist}
\end{outerlist}

\blankline

\href{http://www.filestack.com/}{\textbf{Filepicker.io}}, San Francisco, CA \hfill{Principal Engineer, January 2013  | January 2014}
 \begin{outerlist}
 	\item[] \begin{innerlist}[-]
		\item Scaled a highly available service to \textgreater 300QPS while reducing the
number of servers required
	\end{innerlist}
\end{outerlist}

\blankline

\vspace{-1mm}

\blankline

\section{Skills}
Rust, Java, Scala, Python, Go, (my/Postgre)SQL


\section{Education}
\href{http://mit.edu/}{\textbf{Massachusetts Institute of Technology}},
Cambridge, MA \hfill{June 2013}
\begin{outerlist}
\item[] B.S., \href{http://www.eecs.mit.edu/}{Computer Science and Engineering (6--3)}, \textit{GPA:} 4.9/5.0
\end{outerlist}


\section{Interests}
rock climbing, mountain biking, trail running

\end{document}

%%%%%%%%%%%%%%%%%%%%%%%%%% End CV Document %%%%%%%%%%%%%%%%%%%%%%%%%%%%%
